\section{Certificate tightness and repeated-query economics}\label{sec:r24-diagnostics}
These diagnostics answer different questions from the geometric expansion
theorems. They do not pool new timing observations with historical runs or
reinterpret a fixed-model confidence statement as a training-algorithm
guarantee. The complete measurements and computation protocol are in the
Computational Supplement, Section~\ref{cs:r24-details}.

\subsection{How conservative is the robust economic certificate?}
For each of the 96 recorded contexts and four uncertainty radii, we draw two
strictly interior coefficient vectors from the declared rational grid. This
gives 768 interior models. For each model, we solve both the true accepted
time-only comparator and the common outer comparator at that same objective.
Separate rational feasibility and Fenchel calculations certify lower and upper
values $L_K\leq K_\zeta\leq U_K$ and
$L_O\leq O_\zeta:=\max_{Y^{\rm out}}F_\zeta\leq U_O$.
For the vertex-mixture upper certificate
$U_{\rm mix}=\sum_v\alpha_vU_v$, the actual certificate slack satisfies
\[
 U_{\rm mix}-K_\zeta\in[U_{\rm mix}-U_K,U_{\rm mix}-L_K].
\]
It has two distinct sources:
\[
 U_{\rm mix}-K_\zeta=(O_\zeta-K_\zeta)+(U_{\rm mix}-O_\zeta).
\]
The first is the outer-set relaxation and the second includes vertex
interpolation and the vertex solver's upper-bound error. Both components are
nonnegative. Their computed intervals, rather than an uncertified numerical
optimum, determine the reported summaries.

All 1,536 comparator solves completed. The largest independent comparator
bracket width is \RTwentyFourInteriorGap. Table~\ref{tab:r24-tightness} reports
upper endpoints; the distinction from lower endpoints is at most twice that numerical
scale. At radius one quarter, mean total slack is
\RTwentyFourSlackTotal, split into \RTwentyFourSlackOuter\ from the outer set
and \RTwentyFourSlackInterpolation\ from interpolation. Thus the certificate
is informative but not exact. An interior diagnostic is not a replacement for
the uniform vertex proof and is not a sampled proof of uncertainty-set coverage.

\subsection{Repeated computation against the strongest classical baseline}
The new matched study uses 63-, 127-, and 255-node polyhedra, one separately
fitted 64-label model pair per configuration, 32 held-out contexts, and three
repetitions. The middle configuration has twelve resources and eight extra
capacities; the others have six and four. All methods face the same two exact
certificate targets. Initial method order is randomized within each context;
each tolerance has a separate continuation state. The original eight-context
scaling table remains unchanged, with medians, interquartile ranges, and paired
differences now provided for all its entries in the Computational Supplement.

Every timed decision records prediction, initial optimization, rational audit,
fixed-primal polishing, and fallback setup, solve, and final audit. The fallback
is the same fresh lifted full-objective solve from a zero state for every
method. Therefore no predictor state is credited with saving refinement work
\emph{conditional on fallback}; its possible benefit in this protocol is
avoiding fallback. This distinction prevents warm-start value from being
attributed to a predictor that never initializes the fallback. The full study
contains \RTwentyFourTimingRows\ attempted pipelines and
\RTwentyFourTimingFailures\ failed pipelines; failures, if any, remain in the
raw timing record rather than being discarded.

Let $C_{\rm off}$ be the method's label-generation, fit, and incremental setup
cost relative to lifted continuation, and let $\bar t_L,\bar t_C$ be the
matched per-query mean times. The measured cost difference after $Q$ queries is
\[
 C_{\rm off}+Q(\bar t_L-\bar t_C).
\]
A finite observed crossing exists only when $\bar t_C>\bar t_L$; its threshold
is $C_{\rm off}/(\bar t_C-\bar t_L)$. An interval spanning zero online savings
does not support a stable break-even conclusion. Table~\ref{tab:r24-timing}
reports paired context-cluster bootstrap intervals after averaging the three
repetitions. A moving-block sensitivity analysis also retains short-range dependence from classical continuation. Neither resampling calculation supplies a distribution-free coverage theorem. They describe these fixed fits and sampled contexts, not
arbitrary retraining, tuning, or changing deployment laws. Serialized predictor
sizes, reachable arrays, sparse matrix storage, label times, and individual fit
times are disclosed separately; sparse input storage is not mislabeled as the
solver's total factorization memory.

\subsection{What is required before a decision is called certified?}
There are three distinct protocols. An accepted-policy implementation computes
a response, repairs it, and audits acceptance and its own full-objective gap.
A \emph{pointwise economic certificate} additionally computes a valid upper
bound for the optimized restricted comparator. Under joint uncertainty this
requires all eight common-outer comparator bounds unless previously valid
bounds can be reused. A \emph{population validation} evaluates a frozen policy
on an independent sample, charging those comparator solves to the validation
experiment; it does not make them free online calculations.

The nominal timing tables measure the first protocol, not an online guarantee
of improvement over the optimized restricted contract. The robust experiment's
pointwise certificate uses the second. The archived robust cost analysis now
reports own-decision time plus the complete eight-comparator batch, charged
once to each independently certified method. Its original timers excluded
ensemble prediction; that omission is explicitly identified, so those archived
sums are not advertised as complete learned-pipeline timings. No robust neural
speed claim is inferred. The numerical comparison concerns the third protocol
when the comparator is computed only for offline evaluation. The positive
frozen-ensemble economic-gain bounds remain conditional on those fixed rules,
not on new fitting or model selection.


A separate fresh-host run closes the prediction-timer omission without mixing
machines. It uses sixteen fresh contexts at zero and one-quarter radius and
the unchanged eight-model ensembles. All ninety-six pointwise certifications
include prediction, a fresh candidate workspace, rational repair, vertex
evaluation, and the full eight-comparator batch. At one-quarter radius, mean
complete times are \RTwentyFourRobustTanh, \RTwentyFourRobustRBF, and
\RTwentyFourRobustClassical\ milliseconds for protected tanh, protected RBF,
and classical robust maximin, respectively. Model loading and geometry setup
are disclosed separately from repeated-query cost. Computational Supplement
Table~\ref{tab:r24-robust-cost} reports both candidate and comparator work.
These are descriptive fixed-implementation costs, not a robustness-coverage
or algorithm-level superiority claim.

The robust candidates need not attain the same maximin certificate value. These fresh times are complete cost accounting, not a speed comparison at matched robust gain.
